\section{Open Questions}

The following five questions are genuinely open at the end of this replication, with
concrete next-step probes for each.

\subsection*{Q1. Does the 10$\times$ headline hold beyond H$_2$?}
\textbf{Basis.} This replication tested only H$_2$/4q/STO-3G. The paper's 10$\times$
headline is stated on BeH$_2$ (70$\to$9 CNOTs) and LiH (12 CNOTs), where the HEA baseline
itself needs 60--70+ CNOTs. On H$_2$ the HEA baseline is 6 CNOTs so the maximum possible
reduction factor is small. What we tested is ``sign and mechanism,'' not ``10$\times$ magnitude.''
\textbf{Next steps.} (1) Extend the code to accept a molecule spec. (2) Build BeH$_2$
(linear, 4-electron active, 6 orbitals) and LiH (STO-3G, frozen 1s Li) Hamiltonians via
\texttt{pennylane.qchem.molecular\_hamiltonian}. (3) Scale NSGA-II to pop 32--64, generations
25--50, max CNOT 30--80. (4) Compare to paper's Fig.3--5 point-by-point.
Runtime: $\sim$4--8\,h on uicgpu A100 or $\sim$12--24\,h on m1 CPU.

\subsection*{Q2. Is CMA-ES inner loop fungible with COBYLA/L-BFGS-B on high-D angle spaces?}
\textbf{Basis.} We substituted COBYLA / L-BFGS-B multistart as same-class stand-in. On H$_2$
both hit machine precision. CMA-ES's covariance-matrix adaptation is designed for high-D
noisy landscapes --- which is where the paper actually needs it (12q LiH has 60+ angles).
\textbf{Next steps.} (1) \texttt{pip install cma}, wrap inner objective in a CMA-ES driver
at $\sigma_0=0.5$, maxiter 500. (2) Rerun H$_2$ 3-CNOT topologies with CMA-ES, compare
success rate (currently 7/60 with L-BFGS-B). (3) Head-to-head on the 4-block case. (4) Then
BeH$_2$ as the acid test.

\subsection*{Q3. Is MoG-VQE strictly better than single-objective evolutionary search on the same topology space?}
\textbf{Basis.} The paper's ``multiobjective is necessary'' claim implies single-objective
GA (energy only, no CNOT pressure) would drift to high-CNOT high-parameter circuits and miss
the low-CNOT corner. Our replication compared MoG-VQE against fixed-topology baselines
(UCCSD, HEA), not against a single-objective evolutionary search on the same topology space.
This is the strict test of the ``multi is necessary'' claim.
\textbf{Next steps.} (1) Fork \texttt{code/mog\_vqe\_h2.py} to a single-objective variant
(NSGA-II collapses to standard GA on energy only). (2) Run 5 seeds $\times$ 20 generations
$\times$ pop 32 on H$_2$. (3) Report the final population's CNOT distribution and check
whether the 3-CNOT elbow appears. (4) Add a weighted-sum variant (energy + $\lambda\cdot$NCNOT)
as an intermediate case. \emph{Predicted outcome:} single-obj without penalty concentrates
at high NCNOT (the paper's implicit claim).

\subsection*{Q4. Do modern MOO algorithms (NSGA-III, MOEA/D, SMS-EMOA) beat NSGA-II on this problem?}
\textbf{Basis.} NSGA-II is 20+ years old. NSGA-III (many-objective), MOEA/D
(decomposition-based), SMS-EMOA (hypervolume-driven) may perform better on the deceptive
fitness landscape of quantum-circuit search --- particularly if we add a third objective
(circuit depth or coherent-error metric) for noise-aware MoG-VQE.
\textbf{Next steps.} (1) \texttt{pip install pymoo} (ships all four algorithms). (2) Wrap
topology-genome + inner-optimizer as a pymoo Problem. (3) Run all four with matched
compute budget on H$_2$ and BeH$_2$. (4) Compare hypervolume of final Pareto fronts.
(5) Extension: add third objective = circuit depth $\Rightarrow$ NSGA-III becomes natural.

\subsection*{Q5. Do Pareto fronts transfer across molecules (ML-boosted GA warm-start)?}
\textbf{Basis.} The paper treats each molecule independently, but the family of good
low-CNOT topologies may transfer (e.g., locality of entanglement pattern). Warm-starting
from a related molecule's Pareto front might substantially accelerate convergence, and
intersects the emerging ``foundation models for quantum ansatze'' direction (2024+ ML-VQE).
\textbf{Next steps.} (1) Solve H$_2$ (4q) and H$_4$ (linear chain, 8q) MoG-VQE independently,
save Pareto fronts. (2) Design a topology-lift operator: map $k$-block H$_2$ topology into
8-qubit space by qubit re-mapping + block-doubling. (3) Seed BeH$_2$ NSGA-II 50\% cold +
50\% lifted from H$_4$. (4) Compare generations-to-first-chem-acc vs cold-start.
(5) Related: train a small GNN (SchNet-style on qubit-interaction graph) to predict
which \texttt{(ctrl,tgt)} insertions improve energy, use as ML-boosted mutation operator.
